\documentclass[tikz,border=6pt]{standalone}
\usepackage{peabra}
\begin{document}
\begin{tikzpicture}[x=1cm,y=1cm,
  v/.style={circle, draw=peKek, fill=peHalvany, line width=0.8pt, minimum size=8mm, font=\footnotesize\bfseries},
  t/.style={draw=peSotet, fill=#1, minimum size=6.5mm, font=\footnotesize\bfseries, text=white},
  hi/.style={-{Stealth[length=2mm]}, line width=0.8pt, draw=pePiros},
  lo/.style={-{Stealth[length=2mm]}, line width=0.8pt, draw=peZold, dash pattern=on 3pt off 2pt}]
  \node[v] (x1) at (0,0) {$x_1$};
  \node[v] (g2) at (-1.6,-1.5) {$x_2$};
  \node[v] (h2) at (1.6,-1.5) {$x_2$};
  \node[v] (x4) at (-2.4,-3.0) {$x_4$};
  \node[v] (x3) at (1.6,-3.0) {$x_3$};
  \node[t=pePiros] (one) at (-0.4,-4.5) {1};
  \node[t=peZold] (zero) at (2.0,-4.5) {0};
  \draw[hi] (x1) -- (g2); \draw[lo] (x1) -- (h2);
  \draw[hi] (g2) -- (one); \draw[lo] (g2) -- (x4);
  \draw[hi] (x4) -- (one); \draw[lo] (x4) .. controls (-2.5,-5.3) and (1.2,-5.4) .. (zero);
  \draw[hi] (h2) -- (x3); \draw[lo] (h2) to[bend left=20] (zero);
  \draw[hi] (x3) -- (one); \draw[lo] (x3) -- (zero);
  \node[font=\scriptsize, align=left, anchor=west] at (3.3,-0.4) {\textcolor{pePiros}{\rule[0.6mm]{6mm}{0.8pt}} az esemény bekövetkezik ($x_i=1$)\\[2pt]\textcolor{peZold}{- - -} nem következik be ($x_i=0$)};
  \node[font=\scriptsize, align=left, anchor=west, text width=62mm] at (3.3,-2.4) {$\mathrm{TOP}=(x_1+x_2x_3)(x_2+x_4)$\\[3pt]
    Shannon-kifejtés $x_1$ szerint:\\ $\mathrm{TOP}=x_1(x_2+x_4)+\bar{x}_1\,x_2x_3$\\[3pt]
    a pontos valószínűség az 1-be vezető utak összege:\\ $Q=q_1[q_2+(1-q_2)q_4]+(1-q_1)\,q_2q_3$};
\end{tikzpicture}
\end{document}
